\chapter{Planar rigid transformations and homogeneous coordinates}
\label{chap:lecture03}

Rotation describes the relative orientation of two frames, but a mobile robot
also changes position. Its origin generally differs from the origin of the
world map. A sensor mounted on the robot may have yet another origin and
orientation. Relating point coordinates between these frames requires both
rotation and translation.

This chapter extends the rotation model of Chapter~\ref{chap:lecture02} to
planar rigid transformations. The geometric construction leads first to an
affine coordinate relationship, then to a homogeneous matrix representation
that combines rotation and translation in a single multiplication. The same
representation makes composition and inversion systematic.

\section{Rotation and translation between two frames}
Let $O_W$ and $O_R$ be the origins of world frame $W$ and robot frame $R$.
Define
\[
 p_R^W=\begin{bmatrix}p_{Rx}^W\\p_{Ry}^W\end{bmatrix}
\]
as the position of $O_R$ relative to $O_W$, expressed in world coordinates.
The rotation matrix $R_R^W$ describes the orientation of $R$ relative to $W$.
Together, $p_R^W$ and $R_R^W$ specify the robot frame's planar pose.

For a point $B$, the geometric displacement vectors obey
\begin{equation}
 \overrightarrow{O_WB}=\overrightarrow{O_WO_R}+\overrightarrow{O_RB}.
 \label{eq:l03-vector-sum}
\end{equation}
This triangle relation is independent of coordinates. To turn it into a
numerical calculation, however, every vector being added must be expressed
in the same frame.

\begin{figure}[htbp]
\centering\begin{tikzpicture}[scale=1.35]
\coordinate (O) at (0,0);
\coordinate (R) at (2,1);
\coordinate (B) at (2.366,2.366);
\draw[axis] (O)--(4,0) node[right] {$x_W$};
\draw[axis] (O)--(0,3.2) node[above] {$y_W$};
\foreach \x in {1,2,3} \draw (\x,.04)--(\x,-.04) node[below] {$\x$};
\foreach \y in {1,2} \draw (.04,\y)--(-.04,\y) node[left] {$\y$};
\draw[projection] (B)--(2.366,0);
\draw[projection] (B)--(0,2.366);
\draw[axis,sensorframe] (O)--(R) node[midway,below right] {$p_R^W$};
\draw[axis,thick] (O)--(B) node[pos=.4,above left] {$B^W$};
\draw[axis,sensorframe] (R)--(B) node[midway,right] {$R_R^W B^R$};
\begin{scope}[shift={(R)}]
\draw[axis,robotframe] (0,0)--(30:1.65) node[right] {$x_R$};
\draw[axis,robotframe] (0,0)--(120:1.65) node[above] {$y_R$};
\draw[projection,gray] (0,0)--(1.4,0);
\draw[-{Stealth}] (.65,0) arc (0:30:.65);
\node at (15:.88) {$\theta$};
\end{scope}
\fill (R) circle (1.5pt) node[below] {$O_R$};
\fill (B) circle (1.5pt) node[above right] {$B$};
\node[below left] at (O) {$O_W$};
\end{tikzpicture}

\caption{The displacement to $B$ is the sum of the displacement to the robot
origin and the displacement from that origin to $B$. All three arrow labels
give components in $W$. The geometry uses $p_R^W=[2\;1]\trans$,
$\theta=\pi/6$, and $B^R=[1\;1]\trans$.}
\label{fig:l03-geometry}
\end{figure}

The column $B^R$ contains the components of $\overrightarrow{O_RB}$ along
the robot axes. The column $p_R^W$ uses the world axes. Adding these columns
directly would combine components measured along different directions.
First rotate $B^R$ into world-aligned components; then add the translation:
\begin{equation}
 \boxed{B^W=p_R^W+R_R^W B^R.}
 \label{eq:l03-affine}
\end{equation}
Here $R_R^W B^R$ describes the displacement from $O_R$ to $B$ using world
directions. It is not, by itself, the position of $B$ relative to $O_W$.

\subsection{An intermediate frame makes the geometry explicit}
Introduce a frame $A$ with origin $O_A=O_R$ and axes parallel to those of $W$.
Frames $R$ and $A$ differ only by rotation, so the previous chapter gives
\[
 B^A=R_R^A B^R.
\]
Since $A$ and $W$ have the same orientation, $R_R^A=R_R^W$. Moreover, the
components in $A$ use exactly the directions needed for addition to $p_R^W$.
Thus $B^W=p_R^W+B^A$, recovering Equation~\eqref{eq:l03-affine}.

\begin{figure}[htbp]
\centering\begin{tikzpicture}[scale=1.4]
\coordinate (B) at (.366,1.366);
\draw[axis,sensorframe] (0,0)--(2,0) node[right] {$x_A\parallel x_W$};
\draw[axis,sensorframe] (0,0)--(0,2) node[above] {$y_A\parallel y_W$};
\draw[axis,robotframe] (0,0)--(30:2) node[right] {$x_R$};
\draw[axis,robotframe] (0,0)--(120:1.7) node[above] {$y_R$};
\draw[axis] (0,0)--(B);
\fill (B) circle (1.5pt) node[above right] {$B$};
\draw[projection,sensorframe] (B)--(.366,0) node[below] {$x_B^A$};
\draw[projection,sensorframe] (B)--(0,1.366) node[left] {$y_B^A$};
\draw[projection,robotframe] (B)--(30:1);
\draw[projection,robotframe] (B)--(120:1);
\draw[-{Stealth}] (.7,0) arc (0:30:.7);
\node at (15:.94) {$\theta$};
\node[below left] at (0,0) {$O_A=O_R$};
\end{tikzpicture}

\caption{The auxiliary frame $A$ shares the robot origin but is aligned with
the world. Rotating from $R$ to $A$ resolves the displacement into components
that can be added to the world position of $O_R$.}
\label{fig:l03-intermediate}
\end{figure}

\subsection{Why the relationship is affine}
A map $f(b)=Rb+p$ is called affine: it consists of a linear map followed by
an offset. With a nonzero translation, it is not a linear map of the original
two coordinates, since $f(0)=p\ne0$. In the present setting $R$ is a rotation,
so the transformation is rigid: it preserves distances between points.
Distances to the coordinate origin need not be preserved when the origin changes.

\subsection{Worked example}
Suppose the given quantities are
\[
 B^R=\begin{bmatrix}1\\1\end{bmatrix}\,\mathrm{m},\qquad
 p_R^W=\begin{bmatrix}2\\1\end{bmatrix}\,\mathrm{m},\qquad
 \theta_R^W=\frac{\pi}{6}.
\]
The rotation matrix is dimensionless. Applying the affine relationship gives
\begin{align}
 B^W
 &=\begin{bmatrix}2\\1\end{bmatrix}
 +\begin{bmatrix}\sqrt3/2&-1/2\\1/2&\sqrt3/2\end{bmatrix}
 \begin{bmatrix}1\\1\end{bmatrix}\\
 &=\begin{bmatrix}(3+\sqrt3)/2\\(3+\sqrt3)/2\end{bmatrix}
 \approx\begin{bmatrix}2.3660\\2.3660\end{bmatrix}\,\mathrm{m}.
 \label{eq:l03-example}
\end{align}
The coordinate arithmetic uses metres throughout. The result agrees with
Figure~\ref{fig:l03-geometry}: the point lies above and slightly to the right
of the robot origin. Keeping the exact expressions until the final step avoids
rounding errors when reversing the calculation.

\section{Reversing the affine relationship}
Suppose $B^W$ is known and $B^R$ is required. Subtract the world-coordinate
translation before applying the inverse rotation:
\begin{align}
 B^W-p_R^W&=R_R^W B^R,\\
 \boxed{B^R=(R_R^W)\trans\bigl(B^W-p_R^W\bigr).}
 \label{eq:l03-affine-inverse}
\end{align}
The difference in parentheses is expressed in $W$; the transpose maps it to
$R$. The parentheses matter: the rotation acts on both terms, not just on $B^W$.
Equivalently,
\begin{equation}
 B^R=-(R_R^W)\trans p_R^W+(R_R^W)\trans B^W.
 \label{eq:l03-inverse-expanded}
\end{equation}
This has the same affine structure as the forward relationship, with
\begin{equation}
 R_W^R=(R_R^W)\trans,\qquad p_W^R=-(R_R^W)\trans p_R^W.
 \label{eq:l03-reverse-pose}
\end{equation}
The reversed translation is not generally just $-p_R^W$: it must also be
expressed along the axes of the reverse transformation's output frame.

For the exact coordinates in Equation~\eqref{eq:l03-example},
\[
 B^W-p_R^W=\begin{bmatrix}(\sqrt3-1)/2\\(1+\sqrt3)/2\end{bmatrix}\,\mathrm{m}.
\]
Multiplying by
$\begin{bmatrix}\sqrt3/2&1/2\\-1/2&\sqrt3/2\end{bmatrix}$
recovers $B^R=[1\;1]\trans\,\mathrm{m}$. If the rounded value
$[2.37\;2.37]\trans$ is used instead, the recovered result is approximately
$[1.0054\;1.0015]\trans\,\mathrm{m}$, rather than exactly $[1\;1]\trans$.

\section{Composing transformations through a sensor frame}
A sensor frame $S$ may be displaced and rotated relative to the robot frame.
Given the sensor's pose in $R$ and the robot's pose in $W$, the two coordinate
relationships are
\[
 B^R=p_S^R+R_S^R B^S,\qquad B^W=p_R^W+R_R^W B^R.
\]
Substitution yields
\begin{align}
 B^W&=p_R^W+R_R^W\bigl(p_S^R+R_S^R B^S\bigr)\\
 &=\underbrace{p_R^W+R_R^W p_S^R}_{p_S^W}
 +\underbrace{R_R^W R_S^R}_{R_S^W}B^S.
 \label{eq:l03-composition-expanded}
\end{align}
The overall rotation and translation are therefore
\begin{equation}
 \boxed{R_S^W=R_R^W R_S^R,\qquad
 p_S^W=p_R^W+R_R^W p_S^R.}
\end{equation}
The sensor translation must be rotated into $W$ before it can be added to
the robot translation. Frame labels explain why the expression is valid and
why simply adding $p_R^W+p_S^R$ would be incorrect in general.

\begin{figure}[htbp]
\centering\begin{tikzpicture}[scale=1.1]
\coordinate (R) at (2,.8);
\coordinate (S) at (3.1,2.3);
\coordinate (B) at (4.8,3.4);
\draw[axis] (0,0)--(5.8,0) node[right] {$x_W$};
\draw[axis] (0,0)--(0,4) node[above] {$y_W$};
\node[below left] at (0,0) {$O_W$};
\begin{scope}[shift={(R)}]
\draw[axis,robotframe] (0,0)--(20:1.2) node[right] {$x_R$};
\draw[axis,robotframe] (0,0)--(110:1.2) node[above] {$y_R$};
\end{scope}
\begin{scope}[shift={(S)}]
\draw[axis,sensorframe] (0,0)--(-10:1.2) node[right] {$x_S$};
\draw[axis,sensorframe] (0,0)--(80:1.2) node[above] {$y_S$};
\end{scope}
\draw[axis,black!75] (0,0)--(R) node[midway,below] {$p_R^W$};
\draw[axis,black!75] (R)--(S) node[midway,right] {$R_R^W p_S^R$};
\draw[axis,black!75] (S)--(B) node[midway,below right] {$R_S^W B^S$};
\draw[projection] (0,0)--(B);
\fill (B) circle (1.5pt) node[right] {$B$};
\fill (R) circle (1.5pt) node[below right] {$O_R$};
\fill (S) circle (1.5pt) node[below left] {$O_S$};
\end{tikzpicture}

\caption{A world--robot--sensor chain with distinct origins. Each displacement
along the path to $B$ is labeled in world components so that the arrows can be
added numerically. The dashed line is the direct displacement to $B$.}
\label{fig:l03-composition}
\end{figure}

Repeated substitution works for any number of frames, but each additional
frame introduces another translation term. A representation that absorbs
these additions into matrix multiplication makes longer chains easier to manage.

\section{Homogeneous coordinates and transformation matrices}
\subsection{Adding one coordinate}
For a point with coordinates $B^R$, define its homogeneous representation by
appending a dimensionless one:
\begin{equation}
 \widetilde B^R=\begin{bmatrix}B^R\\1\end{bmatrix}
 =\begin{bmatrix}x_B^R\\y_B^R\\1\end{bmatrix}.
\end{equation}
The additional coordinate is bookkeeping for translation, not a third spatial
coordinate. Define the planar homogeneous transformation matrix as
\begin{equation}
 \boxed{H_R^W=\begin{bmatrix}R_R^W&p_R^W\\0_{1\times2}&1\end{bmatrix}.}
 \label{eq:l03-H}
\end{equation}
Its blocks have dimensions $2\times2$, $2\times1$, $1\times2$, and $1\times1$,
respectively. Written entry by entry, it is
\[
 H_R^W=\begin{bmatrix}
 \cos\theta&-\sin\theta&p_{Rx}^W\\
 \sin\theta&\cos\theta&p_{Ry}^W\\
 0&0&1
 \end{bmatrix}.
\]
The upper-left block contains dimensionless rotation coefficients; the first
two entries of the last column carry the chosen length unit. Block divisions
are visual aids, not extra mathematical operations.

\subsection{One product performs the full transformation}
Block multiplication gives
\begin{equation}
 H_R^W\widetilde B^R
 =\begin{bmatrix}R_R^W&p_R^W\\0_{1\times2}&1\end{bmatrix}
 \begin{bmatrix}B^R\\1\end{bmatrix}
 =\begin{bmatrix}R_R^W B^R+p_R^W\\1\end{bmatrix}
 =\widetilde B^W.
\end{equation}
Thus the affine transformation becomes a single matrix multiplication:
\begin{equation}
 \boxed{\widetilde B^W=H_R^W\widetilde B^R.}
 \label{eq:l03-homogeneous-map}
\end{equation}
The first two output entries are the desired point coordinates. For the
normalized point representation used here, the last entry remains one.
The fixed bottom row $[0\;0\;1]$ is what preserves this property.

For the running example, using metres for the position entries,
\[
 \begin{bmatrix}\sqrt3/2&-1/2&2\\1/2&\sqrt3/2&1\\0&0&1\end{bmatrix}
 \begin{bmatrix}1\\1\\1\end{bmatrix}
 =\begin{bmatrix}(3+\sqrt3)/2\\(3+\sqrt3)/2\\1\end{bmatrix}
 \approx\begin{bmatrix}2.3660\\2.3660\\1\end{bmatrix}.
\]
The physical point coordinates are the first two entries in metres; the final
one has no length unit.

\section{Composition by matrix multiplication}
Using homogeneous coordinates for the sensor chain,
\[
 \widetilde B^R=H_S^R\widetilde B^S,\qquad
 \widetilde B^W=H_R^W\widetilde B^R.
\]
Combining the two maps gives
\begin{equation}
 \boxed{H_S^W=H_R^W H_S^R,\qquad
 \widetilde B^W=H_R^W H_S^R\widetilde B^S.}
\end{equation}
Multiplying the blocks verifies that this contains both composition rules:
\begin{align}
 H_R^W H_S^R
 &=\begin{bmatrix}R_R^W&p_R^W\\0&1\end{bmatrix}
   \begin{bmatrix}R_S^R&p_S^R\\0&1\end{bmatrix}\\
 &=\begin{bmatrix}
 R_R^W R_S^R&p_R^W+R_R^W p_S^R\\0&1
 \end{bmatrix}=H_S^W.
\end{align}
Here each $0$ in the bottom-left block denotes a $1\times2$ zero row.
The product is another homogeneous rigid transformation with bottom row
$[0\;0\;1]$. As before, the rightmost map acts first and adjacent frame
indices must match. For example,
\[
 H_1^2 H_3^1 H_4^3 H_5^4=H_5^2
\]
maps coordinates through $5\to4\to3\to1\to2$.

Unlike pure planar rotations, planar transformations with translation do not
generally commute. The term $R_R^W p_S^R$ shows why: a preceding rotation
changes the direction in which a subsequent local translation acts.

\section{Inverting a homogeneous transformation}
\subsection{Why the transpose is insufficient}
Expanding the determinant along the bottom row gives
\[
 \det(H_R^W)=\det(R_R^W)=1.
\]
Every matrix of this rigid-transformation form is therefore invertible.
However, it is generally not orthogonal. If its translation is nonzero, its
transpose has bottom row $[(p_R^W)\trans\;1]$ rather than $[0\;0\;1]$.
Consequently, transposing the full matrix does not produce its inverse.
When the translation is zero, the full matrix is orthogonal and the transpose
does coincide with the inverse.

\subsection{Deriving the inverse from the affine relationship}
Equation~\eqref{eq:l03-inverse-expanded} already identifies the reverse
rotation and translation. Place these in their corresponding blocks:
\begin{equation}
 \boxed{H_W^R=(H_R^W)^{-1}
 =\begin{bmatrix}
 (R_R^W)\trans&-(R_R^W)\trans p_R^W\\0_{1\times2}&1
 \end{bmatrix}.}
 \label{eq:l03-H-inverse}
\end{equation}
This gives $\widetilde B^R=H_W^R\widetilde B^W$ without a general-purpose
$3\times3$ inversion. Transpose the rotation block, compute the rotated
negative translation, and retain the bottom row.

\begin{figure}[htbp]
\centering\begin{tikzpicture}
\fill[robotframe!10] (0,1) rectangle (2,3);
\fill[sensorframe!12] (2,1) rectangle (3,3);
\draw[thick] (0,0) rectangle (3,3);
\draw[dashed] (2,0)--(2,3) (0,1)--(3,1);
\node at (1,2) {$R_R^W$};
\node at (2.5,2) {$p_R^W$};
\node at (1,.5) {$0\quad0$};
\node at (2.5,.5) {$1$};
\node[above] at (1.5,3.1) {$H_R^W$};
\draw[axis] (3.5,1.5)--(4.7,1.5) node[midway,above] {inverse};
\begin{scope}[shift={(5.2,0)}]
\fill[robotframe!10] (0,1) rectangle (2,3);
\fill[sensorframe!12] (2,1) rectangle (4.2,3);
\draw[thick] (0,0) rectangle (4.2,3);
\draw[dashed] (2,0)--(2,3) (0,1)--(4.2,1);
\node at (1,2) {$(R_R^W)\trans$};
\node at (3.1,2) {$-(R_R^W)\trans p_R^W$};
\node at (1,.5) {$0\quad0$};
\node at (3.1,.5) {$1$};
\node[above] at (2.1,3.1) {$H_W^R$};
\end{scope}
\end{tikzpicture}

\caption{Block structure of a transform and its inverse. Red blocks contain
rotation; blue blocks contain translation. The inverse translation must be
rotated as well as negated. Block widths are schematic.}
\label{fig:l03-blocks}
\end{figure}

For a compact verification, write $R=R_R^W$ and $p=p_R^W$. Then
\[
 \begin{bmatrix}R&p\\0&1\end{bmatrix}
 \begin{bmatrix}R\trans&-R\trans p\\0&1\end{bmatrix}
 =\begin{bmatrix}RR\trans&-RR\trans p+p\\0&1\end{bmatrix}=I_3.
\]
The reverse product also equals $I_3$, since its translation block is
$R\trans p-R\trans p=0$.

\subsection{Worked inverse}
For $p_R^W=[2\;1]\trans\,\mathrm{m}$ and $\theta_R^W=\pi/6$,
\begin{align*}
 -(R_R^W)\trans p_R^W
 &=-\begin{bmatrix}\sqrt3/2&1/2\\-1/2&\sqrt3/2\end{bmatrix}
 \begin{bmatrix}2\\1\end{bmatrix}\\
 &=\begin{bmatrix}-\sqrt3-1/2\\1-\sqrt3/2\end{bmatrix}
 \approx\begin{bmatrix}-2.2321\\0.1340\end{bmatrix}\,\mathrm{m}.
\end{align*}
Hence, with translation entries in metres,
\begin{equation}
 H_W^R=\begin{bmatrix}
 \sqrt3/2&1/2&-\sqrt3-1/2\\
 -1/2&\sqrt3/2&1-\sqrt3/2\\
 0&0&1
 \end{bmatrix}.
\end{equation}
Its last column locates the world origin in robot coordinates. Its upper-left
block describes the world axes in the robot frame. These meanings provide
geometric checks on the inverse, in addition to verifying the matrix product.

\section{Key identities}
\begin{center}
\renewcommand{\arraystretch}{1.35}
\begin{tabular}{ll}
\toprule
Operation & Identity \\
\midrule
Forward point transformation & $B^W=p_R^W+R_R^W B^R$ \\
Reverse point transformation & $B^R=(R_R^W)\trans(B^W-p_R^W)$ \\
Homogeneous point transformation & $\widetilde B^W=H_R^W\widetilde B^R$ \\
Composition & $H_S^W=H_R^W H_S^R$ \\
Reverse rotation & $R_W^R=(R_R^W)\trans$ \\
Reverse translation & $p_W^R=-(R_R^W)\trans p_R^W$ \\
\bottomrule
\end{tabular}
\end{center}
Always specify the frame of a translation, express displacements in a common
frame before adding them, and distinguish the inverse of a rotation matrix
from the inverse of a full homogeneous transformation.

\clearpage
\section{Additional exercises}
Use metres for all position coordinates and counterclockwise-positive angles.
Show your calculations and keep exact values until the final numerical step.
\begin{enumerate}
\item \textbf{A translated and rotated robot.}
Let $p_R^W=[3\;1]\trans$, $\theta_R^W=\pi/2$, and
$B^R=[2\;-1]\trans$. Sketch the two frames and the point. Compute $B^W$
using the affine relationship and explain why $p_R^W+B^R$ gives a different result.

\item \textbf{Recovering robot coordinates.}
Suppose $p_R^W=[1\;-2]\trans$, $\theta_R^W=-\pi/2$, and
$B^W=[4\;-3]\trans$. Find $B^R$. Substitute it into the forward
relationship to verify your answer. State the frame of every vector being added.

\item \textbf{Building a homogeneous matrix.}
Construct $H_R^W$ for $p_R^W=[2\;0]\trans$ and $\theta_R^W=\pi/4$.
Use it to transform $B^R=[\sqrt2\;0]\trans$. Identify the physical
coordinates in the output and explain the role of its last entry.

\item \textbf{World--robot--sensor composition.}
Let $p_R^W=[2\;1]\trans$, $\theta_R^W=\pi/2$,
$p_S^R=[1\;0]\trans$, and $\theta_S^R=-\pi/2$.
Find $H_S^W=H_R^W H_S^R$ and use it to transform $B^S=[1\;2]\trans$.
Check your result by applying the two affine transformations separately.

\item \textbf{Inversion without a general matrix inverse.}
For
\[
 H_R^W=\begin{bmatrix}0&-1&2\\1&0&3\\0&0&1\end{bmatrix},
\]
compute $H_W^R$ using its rotation and translation blocks. Verify
$H_W^R H_R^W=I_3$. Compare your result with $(H_R^W)\trans$ and identify
why the transpose fails.

\item \textbf{Order of operations.}
Consider the two maps, expressed using a fixed coordinate system,
\[
 T=\begin{bmatrix}1&0&1\\0&1&0\\0&0&1\end{bmatrix},\qquad
 Q=\begin{bmatrix}0&-1&0\\1&0&0\\0&0&1\end{bmatrix}.
\]
Compute $QT$ and $TQ$ and apply each to $[0\;0\;1]\trans$.
Explain geometrically why rotating after translating differs from translating
after rotating.

\item \textbf{Distances between points.}
Two points obey $B^W=p+RB^R$ and $C^W=p+RC^R$, with $R$ orthogonal.
Show that $B^W-C^W=R(B^R-C^R)$ and use this to prove that their distance
is preserved. Why need the distance of a single point from its coordinate
origin not be preserved?
\end{enumerate}
